\judulword{11}{0}{C. PERBANDINGAN DAN ANALISIS KRITIS KEDUA JURNAL}
\teksword{8.5}{0}{Bandingkan kedua artikel. Jika aspek tidak dijelaskan, tuliskan “tidak dijelaskan”.}
\begin{tabelword}{|Y|Y|Y|}
\rowcolor{headerword}
\textbf{\ukuranword{8}Aspek} & \textbf{\ukuranword{8}Jurnal 1} & \textbf{\ukuranword{8}Jurnal 2} \\ \hline
Masalah dan tujuan & ......... & ......... \\ \hline
Metode penelitian & ......... & ......... \\ \hline
Teknik analisis & ......... & ......... \\ \hline
Data dan sumber & ......... & ......... \\ \hline
Sampel/objek/subjek & ......... & ......... \\ \hline
Instrumen/perangkat & ......... & ......... \\ \hline
Desain evaluasi & ......... & ......... \\ \hline
Indikator evaluasi & ......... & ......... \\ \hline
Hasil utama & ......... & ......... \\ \hline
Kelebihan & ......... & ......... \\ \hline
Keterbatasan & ......... & ......... \\ \hline
\end{tabelword}

\judulword{10}{4}{Pertanyaan Analisis Perbandingan}
\pertanyaan{9}{1. Apa persamaan dan perbedaan utama kedua penelitian dari sisi masalah, metode, dan bukti?}
\teksword{8.5}{0}{.........}
\pertanyaan{9}{2. Metode mana yang paling sesuai untuk masalah yang diteliti? Jelaskan berdasarkan isi jurnal.}
\teksword{8.5}{0}{.........}
\pertanyaan{9}{3. Penelitian mana yang temuannya lebih meyakinkan? Pertimbangkan kualitas data, metode, dan bukti.}
\teksword{8.5}{0}{.........}
\pertanyaan{9}{4. Apakah kedua pendekatan dapat digabungkan atau saling melengkapi? Jelaskan.}
\teksword{8.5}{0}{.........}

\judulword{10}{4}{D. RESEARCH GAP}
\pertanyaan{9}{1. Apa masalah atau keterbatasan yang belum diselesaikan oleh kedua penelitian?}
\teksword{8.5}{0}{.........}
\pertanyaan{9}{2. Aspek apa yang masih dapat ditingkatkan dari desain, data, metode analisis, atau evaluasi?}
\teksword{8.5}{0}{.........}
\pertanyaan{9}{3. Metode atau pendekatan alternatif apa yang berpotensi memperbaiki hasil?}
\teksword{8.5}{0}{.........}

